\section{Recovering edges from the anticommutator}\label{sec:edges}

Keep the grading of Section~\ref{sec:grading}, with no restriction on the
positive weight $w$. For an operator $H$ that multiplies by a sign, the
anticommutator $MH+HM$ vanishes between opposite signs and multiplies
entries of $M$ by $2$ or $-2$ between equal signs. The following lemma turns this observation into an edge
count even though an arbitrary weight can change the individual edge
entries.

\begin{lemma}[Recovering edges from the grading]\label{lem:edges}
Let $n\ge1$, $\Omega=\{-1,1\}^n$, and $w:\Omega\to(0,\infty)$.
Let $M$ be the centered grading operator \eqref{eq:grading} constructed
from $w$. For $h:\Omega\to\{-1,1\}$, let $H$ be multiplication by $h$
and set
\[
 \mathcal E_h=\{(x,y)\in\Omega^2:\rho(x,y)=1,\ h(x)=h(y)\}.
\]
Then, counting ordered pairs,
\begin{equation}\label{eq:edge-recovery}
 \abs{\mathcal E_h}\le 2\norm{MH+HM}_{\HS}^2.
\end{equation}
\end{lemma}

\begin{proof}
We compare $M$ with the unweighted grading, whose edge entries have
squared modulus $1/4$. First we show that every edge entry of $M$ has
squared modulus at most $1/4$. We then bound the sum of the deficits by
diagonal energy, which
the anticommutator detects for every choice of signs.

\smallskip\noindent\emph{A bound for each off-diagonal entry.}
Write $C_i=[L,Z_i]=[M,Z_i]$, where $[A,B]=AB-BA$.
Its grade blocks are
\[
 \Pi_sC_i\Pi_r=(s-r)\Pi_sZ_i\Pi_r.
\]
We will prove $\norm{C_i}_{\op}\le1$. The commutator keeps only
adjacent-grade blocks, with opposite signs in the two directions.
We first remove the equal-grade blocks by a norm-nonincreasing average,
then supply those opposite signs, up to a common phase, by unitary
conjugation. Introduce the unitaries
\[
 J=\sum_{k=0}^n(-1)^k\Pi_k,\qquad
 W=\sum_{k=0}^n\mathrm i^k\Pi_k,
 \qquad \mathrm i^2=-1.
\]
Set $O_i=(Z_i-JZ_iJ)/2$. The triangle inequality gives
$\norm{O_i}_{\op}\le(\norm{Z_i}_{\op}+\norm{JZ_iJ}_{\op})/2=1$.
In the $(s,r)$ block, the factor defining $O_i$ is
$(1-(-1)^{s+r})/2$. It is zero on the diagonal and one when
$|s-r|=1$. By \eqref{eq:tridiagonal}, $O_i$ therefore consists precisely
of the blocks of $Z_i$ with
$\abs{s-r}=1$. On these blocks, conjugation by $W$ multiplies by
$\mathrm i^{s-r}=\mathrm i(s-r)$. Therefore
\[
 WO_iW^*=\mathrm i C_i,
 \qquad\text{and hence}\qquad \norm{C_i}_{\op}\le1.
\]
In the point basis, $(C_i)_{xy}=(y_i-x_i)M_{xy}$.
If $x\ne y$, choose $i$ with $x_i\ne y_i$ to obtain
\[
 2\abs{M_{xy}}=\abs{(C_i)_{xy}}\le\norm{C_i}_{\op}\le1.
\]
In particular, each neighbor deficit $1/4-\abs{M_{xy}}^2$ is nonnegative.

\smallskip\noindent\emph{The total neighbor deficit.}
The same grade blocks also give an aggregate bound. By
\eqref{eq:block-parseval} and \eqref{eq:tridiagonal},
\[
 \norm{C_i}_{\HS}^2
 =\sum_{r,s=0}^n(s-r)^2\norm{\Pi_sZ_i\Pi_r}_{\HS}^2
 \le\norm{Z_i}_{\HS}^2=N.
\]
Summing the squared point-basis entries of $C_i$ over coordinates gives
\begin{equation}\label{eq:distance-energy}
 \sum_{x,y}\rho(x,y)\abs{M_{xy}}^2
 =\frac14\sum_{i=1}^n\norm{C_i}_{\HS}^2
 \le\frac{nN}{4}=\sum_{x,y}\abs{M_{xy}}^2,
\end{equation}
where the last equality is \eqref{eq:total-energy}. Set
\[
 A_0=\sum_x\abs{M_{xx}}^2,\qquad
 F=\sum_{\rho(x,y)\ge2}\abs{M_{xy}}^2.
\]
The difference between the outermost sums in
\eqref{eq:distance-energy} is
$-A_0+\sum_{\rho(x,y)\ge2}(\rho(x,y)-1)|M_{xy}|^2\le0$.
Thus energy beyond neighboring pairs is bounded by diagonal energy:
\begin{equation}\label{eq:far-energy}
 F\le\sum_{\rho(x,y)\ge2}(\rho(x,y)-1)\abs{M_{xy}}^2\le A_0.
\end{equation}

There are $nN$ ordered neighbor pairs, and thus
\begin{align}
 \sum_{\rho(x,y)=1}\left(\frac14-\abs{M_{xy}}^2\right)
 &=\frac{nN}{4}-\sum_{\rho(x,y)=1}\abs{M_{xy}}^2\notag\\
 &=A_0+F\le2A_0.\label{eq:deficit}
\end{align}
Here we used \eqref{eq:total-energy} and \eqref{eq:far-energy}.
The fixed total energy consequently controls the full deficit from the
unweighted edge value $1/4$.

\smallskip\noindent\emph{Recovering the equal-sign edges.}
Let $A_h=\sum_{(x,y)\in\mathcal E_h}\abs{M_{xy}}^2$.
Because the individual deficits are nonnegative, their sum over
$\mathcal E_h$ is at most the full sum in \eqref{eq:deficit}. Hence
\begin{equation}\label{eq:count-deficit}
 \frac{\abs{\mathcal E_h}}4\le A_h+2A_0\le2(A_h+A_0).
\end{equation}
For $T=MH+HM$ we have
$T_{xy}=(h(x)+h(y))M_{xy}$. On both the diagonal and $\mathcal E_h$,
the squared factor is $4$, so
\[
 4(A_0+A_h)\le\norm T_{\HS}^2.
\]
Together with \eqref{eq:count-deficit}, this proves
\eqref{eq:edge-recovery}.
\end{proof}
